% arara: pdflatex: {shell: yes, files: [latexindent]}
% the -m switch
% the -m switch
% the -m switch

\fancyhead[R]{\bfseries\thepage%
\tikz[remember picture,overlay] {
\node at (1,0){\includegraphics{logo}};
}}
\section{The \texttt{-m} (\texttt{modifylinebreaks}) switch}\label{sec:modifylinebreaks}
 All features described in this section will only be relevant if the \texttt{-m} switch
 is used.

\yamltitle{modifylinebreaks}*{fields}
	\begin{wrapfigure}[7]{r}[0pt]{8cm}
	\cmhlistingsfromfile[firstnumber=345,linerange={345-347},style=yaml-LST,numbers=left,]{../defaultSettings.yaml}[MLB-TCB,width=.85\linewidth,before=\centering]{\texttt{modifyLineBreaks}}{lst:modifylinebreaks}
\end{wrapfigure}
	\makebox[0pt][r]{%
	\raisebox{-\totalheight}[0pt][0pt]{%
	\tikz\node[opacity=1] at (0,0) {\includegraphics[width=4cm]{logo}};}}%	
	One of the most exciting features of Version 3.0 is the \texttt{-m} switch, which
	permits \texttt{latexindent.pl} to modify line breaks, according to the
	specifications in the \texttt{modifyLineBreaks} field. \emph{The settings
		in this field will only be considered if the \texttt{-m} switch has been used}.
	A snippet of the default settings of this field is shown in \cref{lst:modifylinebreaks}.

	Having read the previous paragraph, it should sound reasonable that, if you call \texttt{latexindent.pl}
	using the \texttt{-m} switch, then you give it permission to modify line breaks in your file,
	but let's be clear:

	\begin{warning}
		If you call \texttt{latexindent.pl} with the \texttt{-m} switch, then you
		are giving it permission to modify line breaks. By default, the only
		thing that will happen is that multiple blank lines will be condensed into
		one blank line; many other settings are possible, discussed next.
	\end{warning}

\yamltitle{preserveBlankLines}{0|1}
	This field is directly related to \emph{poly-switches}, discussed below.
	By default, it is set to \texttt{1}, which means that blank lines will
	be protected from removal; however, regardless of this setting, multiple
	blank lines can be condensed if \texttt{condenseMultipleBlankLinesInto} is
	greater than \texttt{0}, discussed next.

\yamltitle{condenseMultipleBlankLinesInto}*{integer $\geq 0$}
	Assuming that this switch takes an integer value greater than \texttt{0}, \texttt{latexindent.pl} will condense multiple blank lines into
	the number of blank lines illustrated by this switch. As an example, \cref{lst:mlb-bl} shows a sample file
	with blank lines; upon running
	\begin{commandshell}
latexindent.pl myfile.tex -m  
\end{commandshell}
	the output is shown in \cref{lst:mlb-bl-out}; note that the multiple blank lines have been
	condensed into one blank line, and note also that we have used the \texttt{-m} switch!

	\begin{minipage}{.45\textwidth}
		\cmhlistingsfromfile{demonstrations/mlb1.tex}{\texttt{mlb1.tex}}{lst:mlb-bl}
	\end{minipage}%
	\hfill
	\begin{minipage}{.45\textwidth}
		\cmhlistingsfromfile{demonstrations/mlb1-out.tex}{\texttt{mlb1.tex} out output}{lst:mlb-bl-out}
	\end{minipage}

\subsection{Poly-switches}
	Every other field in the \texttt{modifyLineBreaks} field uses poly-switches, and can take
	one of four integer values\footnote{visual learners might like to associate one of the four circles in the logo with one of the four given values}:
	\begin{itemize}[font=\bfseries]
		\item[$-1$] \emph{remove mode}: line breaks before or after the \emph{<part of thing>} can be removed (assuming that \texttt{preserveBlankLines} is set to \texttt{0});
		\item[0] \emph{off mode}: line breaks will not be modified for the \emph{<part of thing>} under consideration;
		\item[1] \emph{add mode}: a line break will be added before or after the \emph{<part of thing>} under consideration, assuming that
		      there is not already a line break before or after the \emph{<part of thing>};
		\item[2] \emph{comment then add mode}: a comment symbol will be added, followed by a line break before or after the \emph{<part of thing>} under consideration, assuming that
		      there is not already a comment and line break before or after the \emph{<part of thing>}.
	\end{itemize}
	All poly-switches are \emph{off} by default; \texttt{latexindent.pl} searches first of all for per-name settings, and then followed by global per-thing settings.

\subsection{modifyLineBreaks for environments}\label{sec:modifylinebreaks-environments}
	We start by viewing a snippet of \texttt{defaultSettings.yaml} in \cref{lst:environments-mlb}; note that it contains \emph{global} settings (immediately
	after the \texttt{environments} field) and that \emph{per-name} settings are also allowed -- in the case of \cref{lst:environments-mlb}, settings
	for \texttt{equation*} have been specified. Note that all poly-switches are \emph{off} by default.

	\cmhlistingsfromfile[firstnumber=348,linerange={348-357},style=yaml-LST,numbers=left,]{../defaultSettings.yaml}[width=.8\linewidth,before=\centering,MLB-TCB]{\texttt{environments}}{lst:environments-mlb}

\subsubsection{Adding line breaks (poly-switches set to $1$ or $2$)}
	Let's begin with the simple example given in \cref{lst:env-mlb1-tex}; note that we have annotated key parts of the file using $\BeginStartsOnOwnLine$,
	$\BodyStartsOnOwnLine$, $\EndStartsOnOwnLine$ and $\EndFinishesWithLineBreak$, these will be related to fields specified in \cref{lst:environments-mlb}.

	\begin{cmhlistings}[escapeinside={(*@}{@*)}]{\texttt{env-mlb1.tex}}{lst:env-mlb1-tex}
before words(*@$\BeginStartsOnOwnLine$@*) \begin{myenv}(*@$\BodyStartsOnOwnLine$@*)body of myenv(*@$\EndStartsOnOwnLine$@*)\end{myenv}(*@$\EndFinishesWithLineBreak$@*) after words
\end{cmhlistings}

	Let's explore \texttt{BeginStartsOnOwnLine} and \texttt{BodyStartsOnOwnLine} in \cref{lst:env-mlb1,lst:env-mlb2}, and in particular,
	let's allow each of them in turn to take a value of $1$.

	\begin{minipage}{.45\textwidth}
		\cmhlistingsfromfile[style=yaml-LST]{demonstrations/env-mlb1.yaml}[MLB-TCB]{\texttt{env-mlb1.yaml}}{lst:env-mlb1}
	\end{minipage}
	\hfill
	\begin{minipage}{.45\textwidth}
		\cmhlistingsfromfile[style=yaml-LST]{demonstrations/env-mlb2.yaml}[MLB-TCB]{\texttt{env-mlb2.yaml}}{lst:env-mlb2}
	\end{minipage}

	After running the following commands,
	\begin{commandshell}
latexindent.pl -m env-mlb.tex -l env-mlb1.yaml
latexindent.pl -m env-mlb.tex -l env-mlb2.yaml
\end{commandshell}
	the output is as in \cref{lst:env-mlb-mod1,lst:env-mlb-mod2} respectively.

	\begin{widepage}
		\begin{minipage}{.57\linewidth}
			\cmhlistingsfromfile{demonstrations/env-mlb-mod1.tex}{\texttt{env-mlb.tex} using \cref{lst:env-mlb1}}{lst:env-mlb-mod1}
		\end{minipage}
		\hfill
		\begin{minipage}{.42\linewidth}
			\cmhlistingsfromfile{demonstrations/env-mlb-mod2.tex}{\texttt{env-mlb.tex} using \cref{lst:env-mlb2}}{lst:env-mlb-mod2}
		\end{minipage}
	\end{widepage}

	There are a couple of points to note:
	\begin{itemize}
		\item in \cref{lst:env-mlb-mod1} a line break has been added at the point denoted by $\BeginStartsOnOwnLine$ in \cref{lst:env-mlb1-tex}; no
		      other line breaks have been changed;
		\item in \cref{lst:env-mlb-mod2} a line break has been added at the point denoted by $\BodyStartsOnOwnLine$ in \cref{lst:env-mlb1-tex};
		      furthermore, note that the \emph{body} of \texttt{myenv} has received the appropriate (default) indentation.
	\end{itemize}

	Let's now change each of the \texttt{1} values in \cref{lst:env-mlb1,lst:env-mlb2} so that they are $2$ and
	save them into \texttt{env-mlb3.yaml} and \texttt{env-mlb4.yaml} respectively (see \cref{lst:env-mlb3,lst:env-mlb4}).

	\begin{minipage}{.45\textwidth}
		\cmhlistingsfromfile[style=yaml-LST]{demonstrations/env-mlb3.yaml}[MLB-TCB]{\texttt{env-mlb3.yaml}}{lst:env-mlb3}
	\end{minipage}
	\hfill
	\begin{minipage}{.45\textwidth}
		\cmhlistingsfromfile[style=yaml-LST]{demonstrations/env-mlb4.yaml}[MLB-TCB]{\texttt{env-mlb4.yaml}}{lst:env-mlb4}
	\end{minipage}

	Upon running  commands analogous to the above, we obtain \cref{lst:env-mlb-mod3,lst:env-mlb-mod4}.

	\begin{widepage}
		\begin{minipage}{.57\linewidth}
			\cmhlistingsfromfile{demonstrations/env-mlb-mod3.tex}{\texttt{env-mlb.tex} using \cref{lst:env-mlb3}}{lst:env-mlb-mod3}
		\end{minipage}
		\hfill
		\begin{minipage}{.42\linewidth}
			\cmhlistingsfromfile{demonstrations/env-mlb-mod4.tex}{\texttt{env-mlb.tex} using \cref{lst:env-mlb4}}{lst:env-mlb-mod4}
		\end{minipage}
	\end{widepage}

	Note that line breaks have been added as in \cref{lst:env-mlb-mod1,lst:env-mlb-mod2}, but this time a comment symbol
	has been added before adding the line break; in both cases, trailing horizontal
	space has been stripped before doing so.

	Let's explore \texttt{EndStartsOnOwnLine} and \texttt{EndFinishesWithLineBreak} in \cref{lst:env-mlb5,lst:env-mlb6},
	and in particular, let's allow each of them in turn to take a value of $1$.

	\begin{minipage}{.49\textwidth}
		\cmhlistingsfromfile[style=yaml-LST]{demonstrations/env-mlb5.yaml}[MLB-TCB]{\texttt{env-mlb5.yaml}}{lst:env-mlb5}
	\end{minipage}
	\hfill
	\begin{minipage}{.49\textwidth}
		\cmhlistingsfromfile[style=yaml-LST]{demonstrations/env-mlb6.yaml}[MLB-TCB]{\texttt{env-mlb6.yaml}}{lst:env-mlb6}
	\end{minipage}

	After running the following commands,
	\begin{commandshell}
latexindent.pl -m env-mlb.tex -l env-mlb5.yaml
latexindent.pl -m env-mlb.tex -l env-mlb6.yaml
\end{commandshell}
	the output is as in \cref{lst:env-mlb-mod5,lst:env-mlb-mod6}.

	\begin{widepage}
		\begin{minipage}{.42\linewidth}
			\cmhlistingsfromfile{demonstrations/env-mlb-mod5.tex}{\texttt{env-mlb.tex} using \cref{lst:env-mlb5}}{lst:env-mlb-mod5}
		\end{minipage}
		\hfill
		\begin{minipage}{.57\linewidth}
			\cmhlistingsfromfile{demonstrations/env-mlb-mod6.tex}{\texttt{env-mlb.tex} using \cref{lst:env-mlb6}}{lst:env-mlb-mod6}
		\end{minipage}
	\end{widepage}

	There are a couple of points to note:
	\begin{itemize}
		\item in \cref{lst:env-mlb-mod5} a line break has been added at the point denoted by $\EndStartsOnOwnLine$ in \vref{lst:env-mlb1-tex}; no
		      other line breaks have been changed and the \lstinline!\end{myenv}! statement has \emph{not} received indentation (as intended);
		\item in \cref{lst:env-mlb-mod6} a line break has been added at the point denoted by $\EndFinishesWithLineBreak$ in \vref{lst:env-mlb1-tex}.
	\end{itemize}

	Let's now change each of the \texttt{1} values in \cref{lst:env-mlb5,lst:env-mlb6} so that they are $2$ and
	save them into \texttt{env-mlb7.yaml} and \texttt{env-mlb8.yaml} respectively (see \cref{lst:env-mlb7,lst:env-mlb8}).

	\begin{minipage}{.49\textwidth}
		\cmhlistingsfromfile[style=yaml-LST]{demonstrations/env-mlb7.yaml}[MLB-TCB]{\texttt{env-mlb7.yaml}}{lst:env-mlb7}
	\end{minipage}
	\hfill
	\begin{minipage}{.49\textwidth}
		\cmhlistingsfromfile[style=yaml-LST]{demonstrations/env-mlb8.yaml}[MLB-TCB]{\texttt{env-mlb8.yaml}}{lst:env-mlb8}
	\end{minipage}

	Upon running  commands analogous to the above, we obtain \cref{lst:env-mlb-mod7,lst:env-mlb-mod8}.

	\begin{widepage}
		\begin{minipage}{.42\linewidth}
			\cmhlistingsfromfile{demonstrations/env-mlb-mod7.tex}{\texttt{env-mlb.tex} using \cref{lst:env-mlb7}}{lst:env-mlb-mod7}
		\end{minipage}
		\hfill
		\begin{minipage}{.57\linewidth}
			\cmhlistingsfromfile{demonstrations/env-mlb-mod8.tex}{\texttt{env-mlb.tex} using \cref{lst:env-mlb8}}{lst:env-mlb-mod8}
		\end{minipage}
	\end{widepage}

	Note that line breaks have been added as in \cref{lst:env-mlb-mod5,lst:env-mlb-mod6}, but this time a comment symbol
	has been added before adding the line break; in both cases, trailing horizontal
	space has been stripped before doing so.

	If you ask \texttt{latexindent.pl} to add a line break (possibly with a comment) using a poly-switch value of $1$ (or $2$),
	it will only do so if necessary. For example, if you process the file in \vref{lst:mlb2} using any of the YAML
	files presented so far in this section, it will be left unchanged.

	\begin{minipage}{.45\linewidth}
		\cmhlistingsfromfile{demonstrations/env-mlb2.tex}{\texttt{env-mlb2.tex}}{lst:mlb2}
	\end{minipage}
	\hfill
	\begin{minipage}{.45\linewidth}
		\cmhlistingsfromfile{demonstrations/env-mlb3.tex}{\texttt{env-mlb3.tex}}{lst:mlb3}
	\end{minipage}

	In contrast, the output from processing the file in \cref{lst:mlb3} will vary depending
	on the poly-switches used; in \cref{lst:env-mlb3-mod2} you'll see that the comment symbol after
	the \lstinline!\begin{myenv}! has been moved to the next line, as \texttt{BodyStartsOnOwnLine}
	is set to \texttt{1}. In \cref{lst:env-mlb3-mod4} you'll see that the comment has been accounted
	for correctly because \texttt{BodyStartsOnOwnLine} has been set to \texttt{2},
	and the comment symbol has \emph{not} been moved to its own line. You're encouraged to experiment
	with \cref{lst:mlb3} and by setting the other poly-switches considered so far to \texttt{2} in turn.

	\begin{minipage}{.45\linewidth}
		\cmhlistingsfromfile{demonstrations/env-mlb3-mod2.tex}{\texttt{env-mlb3.tex} using \vref{lst:env-mlb2}}{lst:env-mlb3-mod2}
	\end{minipage}
	\hfill
	\begin{minipage}{.45\linewidth}
		\cmhlistingsfromfile{demonstrations/env-mlb3-mod4.tex}{\texttt{env-mlb3.tex} using \vref{lst:env-mlb4}}{lst:env-mlb3-mod4}
	\end{minipage}

	The details of the discussion in this section have concerned \emph{global} poly-switches in the \texttt{environments} field;
	each switch can also be specified on a \emph{per-name} basis, which would take priority over the global values; with
	reference to \vref{lst:environments-mlb}, an example is shown for the \texttt{equation*} environment.

\subsubsection{Removing line breaks (poly-switches set to $-1$)}
	Setting poly-switches to $-1$ tells \texttt{latexindent.pl} to remove line breaks of the \emph{<part of the thing>}, if necessary. We will consider the
	example code given in \cref{lst:mlb4}, noting in particular the positions of
	the line break highlighters, $\BeginStartsOnOwnLine$, $\BodyStartsOnOwnLine$, $\EndStartsOnOwnLine$
	and $\EndFinishesWithLineBreak$, together with the associated YAML files in \crefrange{lst:env-mlb9}{lst:env-mlb12}.

	\begin{minipage}{.45\linewidth}
		\begin{cmhlistings}[escapeinside={(*@}{@*)}]{\texttt{env-mlb4.tex}}{lst:mlb4}
before words(*@$\BeginStartsOnOwnLine$@*)
\begin{myenv}(*@$\BodyStartsOnOwnLine$@*)
body of myenv(*@$\EndStartsOnOwnLine$@*)
\end{myenv}(*@$\EndFinishesWithLineBreak$@*)
after words
\end{cmhlistings}
	\end{minipage}%
	\hfill
	\begin{minipage}{.51\textwidth}
		\cmhlistingsfromfile[style=yaml-LST]{demonstrations/env-mlb9.yaml}[MLB-TCB]{\texttt{env-mlb9.yaml}}{lst:env-mlb9}

		\cmhlistingsfromfile[style=yaml-LST]{demonstrations/env-mlb10.yaml}[MLB-TCB]{\texttt{env-mlb10.yaml}}{lst:env-mlb10}

		\cmhlistingsfromfile[style=yaml-LST]{demonstrations/env-mlb11.yaml}[MLB-TCB]{\texttt{env-mlb11.yaml}}{lst:env-mlb11}

		\cmhlistingsfromfile[style=yaml-LST]{demonstrations/env-mlb12.yaml}[MLB-TCB]{\texttt{env-mlb12.yaml}}{lst:env-mlb12}
	\end{minipage}

	After running the commands
	\begin{commandshell}
latexindent.pl -m env-mlb4.tex -l env-mlb9.yaml
latexindent.pl -m env-mlb4.tex -l env-mlb10.yaml
latexindent.pl -m env-mlb4.tex -l env-mlb11.yaml
latexindent.pl -m env-mlb4.tex -l env-mlb12.yaml
\end{commandshell}

	we obtain the respective output in \crefrange{lst:env-mlb4-mod9}{lst:env-mlb4-mod12}.

	\begin{minipage}{.45\linewidth}
		\cmhlistingsfromfile{demonstrations/env-mlb4-mod9.tex}{\texttt{env-mlb4.tex} using \cref{lst:env-mlb9}}{lst:env-mlb4-mod9}
	\end{minipage}
	\hfill
	\begin{minipage}{.45\linewidth}
		\cmhlistingsfromfile{demonstrations/env-mlb4-mod10.tex}{\texttt{env-mlb4.tex} using \cref{lst:env-mlb10}}{lst:env-mlb4-mod10}
	\end{minipage}

	\begin{minipage}{.45\linewidth}
		\cmhlistingsfromfile{demonstrations/env-mlb4-mod11.tex}{\texttt{env-mlb4.tex} using \cref{lst:env-mlb11}}{lst:env-mlb4-mod11}
	\end{minipage}
	\hfill
	\begin{minipage}{.45\linewidth}
		\cmhlistingsfromfile{demonstrations/env-mlb4-mod12.tex}{\texttt{env-mlb4.tex} using \cref{lst:env-mlb12}}{lst:env-mlb4-mod12}
	\end{minipage}

	Notice that in
	\begin{itemize}
		\item \cref{lst:env-mlb4-mod9} the line break denoted by $\BeginStartsOnOwnLine$ in \cref{lst:mlb4} has been removed;
		\item \cref{lst:env-mlb4-mod10} the line break denoted by $\BodyStartsOnOwnLine$ in \cref{lst:mlb4} has been removed;
		\item \cref{lst:env-mlb4-mod11} the line break denoted by $\EndStartsOnOwnLine$ in \cref{lst:mlb4} has been removed;
		\item \cref{lst:env-mlb4-mod12} the line break denoted by $\EndFinishesWithLineBreak$ in \cref{lst:mlb4} has been removed.
	\end{itemize}
	We examined each of these cases separately for clarity of explanation, but you can combine all of the YAML
	settings in \crefrange{lst:env-mlb9}{lst:env-mlb12} into one file; alternatively, you could tell \texttt{latexindent.pl}
	to load them all by using the following command, for example
	\begin{widepage}
		\begin{commandshell}
latexindent.pl -m env-mlb4.tex -l env-mlb9.yaml,env-mlb10.yaml,env-mlb11.yaml,env-mlb12.yaml
\end{commandshell}
	\end{widepage}
	which gives the output in \vref{lst:env-mlb1-tex}.

	\paragraph{About trailing horizontal space}
	Recall that on \cpageref{yaml:removeTrailingWhitespace} we discussed the YAML field \texttt{removeTrailingWhitespace},
	and that it has two (binary) switches to determine if horizontal space should be removed \texttt{beforeProcessing} and \texttt{afterProcessing}.
	The \texttt{beforeProcessing} is particularly relevant when considering the \texttt{-m} switch; let's consider the
	file shown in \cref{lst:mlb5}, which highlights trailing spaces.

	\begin{minipage}{.45\linewidth}
		\begin{cmhlistings}[showspaces=true,escapeinside={(*@}{@*)}]{\texttt{env-mlb5.tex}}{lst:mlb5}
before words   (*@$\BeginStartsOnOwnLine$@*) 
\begin{myenv}           (*@$\BodyStartsOnOwnLine$@*)
body of myenv      (*@$\EndStartsOnOwnLine$@*) 
\end{myenv}     (*@$\EndFinishesWithLineBreak$@*)
after words
\end{cmhlistings}
	\end{minipage}
	\hfill
	\begin{minipage}{.45\linewidth}
		\cmhlistingsfromfile[style=yaml-LST]{demonstrations/removeTWS-before.yaml}[yaml-TCB]{\texttt{removeTWS-before.yaml}}{lst:removeTWS-before}
	\end{minipage}

	The output from the following commands
	\begin{widepage}
		\begin{commandshell}
latexindent.pl -m env-mlb5.tex -l env-mlb9.yaml,env-mlb10.yaml,env-mlb11.yaml,env-mlb12.yaml
latexindent.pl -m env-mlb5.tex -l env-mlb9.yaml,env-mlb10.yaml,env-mlb11.yaml,env-mlb12.yaml,removeTWS-before.yaml
\end{commandshell}
	\end{widepage}
	is shown, respectively, in \cref{lst:env-mlb5-modAll,lst:env-mlb5-modAll-remove-WS}; note that
	the trailing horizontal white space has been preserved (by default) in \cref{lst:env-mlb5-modAll}, while
	in \cref{lst:env-mlb5-modAll-remove-WS}, it has been removed using the switch specified in \cref{lst:removeTWS-before}.

    \begin{widepage}
	\cmhlistingsfromfile{demonstrations/env-mlb5-modAll.tex}{\texttt{env-mlb5.tex} using \crefrange{lst:env-mlb4-mod9}{lst:env-mlb4-mod12}}{lst:env-mlb5-modAll}

	\cmhlistingsfromfile{demonstrations/env-mlb5-modAll-remove-WS.tex}{\texttt{env-mlb5.tex} using \crefrange{lst:env-mlb4-mod9}{lst:env-mlb4-mod12} \emph{and} \cref{lst:removeTWS-before}}{lst:env-mlb5-modAll-remove-WS}
    \end{widepage}


	\paragraph{Blank lines}
	Now let's consider the file in \cref{lst:mlb6}, which contains blank lines.

	\begin{minipage}{.45\linewidth}
		\begin{cmhlistings}[escapeinside={(*@}{@*)}]{\texttt{env-mlb6.tex}}{lst:mlb6}
before words(*@$\BeginStartsOnOwnLine$@*)


\begin{myenv}(*@$\BodyStartsOnOwnLine$@*)


body of myenv(*@$\EndStartsOnOwnLine$@*)


\end{myenv}(*@$\EndFinishesWithLineBreak$@*)

after words
\end{cmhlistings}
	\end{minipage}%
	\hfill
	\begin{minipage}{.45\linewidth}
		\cmhlistingsfromfile[style=yaml-LST]{demonstrations/UnpreserveBlankLines.yaml}[MLB-TCB]{\texttt{UnpreserveBlankLines.yaml}}{lst:UnpreserveBlankLines}
	\end{minipage}

	Upon running the following commands
	\begin{widepage}
		\begin{commandshell}
latexindent.pl -m env-mlb6.tex -l env-mlb9.yaml,env-mlb10.yaml,env-mlb11.yaml,env-mlb12.yaml
latexindent.pl -m env-mlb6.tex -l env-mlb9.yaml,env-mlb10.yaml,env-mlb11.yaml,env-mlb12.yaml,UnpreserveBlankLines.yaml
\end{commandshell}
	\end{widepage}
	we receive the respective outputs in \cref{lst:env-mlb6-modAll,lst:env-mlb6-modAll-un-Preserve-Blank-Lines}. In
	\cref{lst:env-mlb6-modAll} we see that the multiple blank lines have each been condensed into one blank line,
	but that blank lines have \emph{not} been removed by the poly-switches -- this is because, by default, \texttt{preserveBlankLines}
	is set to \texttt{1}. By contrast, in \cref{lst:env-mlb6-modAll-un-Preserve-Blank-Lines}, we have allowed
	the poly-switches to remove blank lines because, in \cref{lst:UnpreserveBlankLines}, we have set \texttt{preserveBlankLines} to \texttt{0}.

	\begin{minipage}{.30\linewidth}
		\cmhlistingsfromfile{demonstrations/env-mlb6-modAll.tex}{\texttt{env-mlb6.tex} using \crefrange{lst:env-mlb4-mod9}{lst:env-mlb4-mod12}}{lst:env-mlb6-modAll}
	\end{minipage}
	\hfill
	\begin{minipage}{.7\linewidth}
		\cmhlistingsfromfile{demonstrations/env-mlb6-modAll-un-Preserve-Blank-Lines.tex}{\texttt{env-mlb6.tex} using \crefrange{lst:env-mlb4-mod9}{lst:env-mlb4-mod12} \emph{and} \cref{lst:UnpreserveBlankLines}}{lst:env-mlb6-modAll-un-Preserve-Blank-Lines}
	\end{minipage}

\subsection{Poly-switches for other code blocks}
	Rather than repeat the examples shown for the environment code blocks (in \vref{sec:modifylinebreaks-environments}), we choose to detail the poly-switches for
	all other code blocks in \cref{tab:poly-switch-mapping}; note that each and every one of these poly-switches is \emph{off by default}.

	\begin{longtable}{m{.2\textwidth}@{\hspace{.75cm}}m{.35\textwidth}@{}m{.4\textwidth}}
		\caption{Poly-switch mappings for all code-block types}\label{tab:poly-switch-mapping}\\
		\toprule
		Code block & Sample & Poly-switch mapping \\
		\midrule
		environment &
		\begin{lstlisting}[escapeinside={(*@}{@*)},nolol=true]
before words(*@$\BeginStartsOnOwnLine$@*)
\begin{myenv}(*@$\BodyStartsOnOwnLine$@*)
body of myenv(*@$\EndStartsOnOwnLine$@*)
\end{myenv}(*@$\EndFinishesWithLineBreak$@*)
after words
  \end{lstlisting}
		&
		\begin{tabular}[t]{c@{~}l@{}}
			$\BeginStartsOnOwnLine$     & BeginStartsOnOwnLine     \\
			$\BodyStartsOnOwnLine$      & BodyStartsOnOwnLine      \\
			$\EndStartsOnOwnLine$       & EndStartsOnOwnLine       \\
			$\EndFinishesWithLineBreak$ & EndFinishesWithLineBreak \\
		\end{tabular}
		\\
		\cmidrule{2-3}
		ifelsefi &
		\begin{lstlisting}[escapeinside={(*@}{@*)},nolol=true]
before words(*@$\BeginStartsOnOwnLine$@*)
\if...(*@$\BodyStartsOnOwnLine$@*)
body of if statement(*@$\ElseStartsOnOwnLine$@*)
\else(*@$\ElseFinishesWithLineBreak$@*)
body of else statement(*@$\EndStartsOnOwnLine$@*)
\fi(*@$\EndFinishesWithLineBreak$@*)
after words
  \end{lstlisting}
		&
		\begin{tabular}[t]{c@{~}l@{}}
			$\BeginStartsOnOwnLine$      & IfStartsOnOwnLine         \\
			$\BodyStartsOnOwnLine$       & BodyStartsOnOwnLine       \\
			$\ElseStartsOnOwnLine$       & ElseStartsOnOwnLine       \\
			$\ElseFinishesWithLineBreak$ & ElseFinishesWithLineBreak \\
			$\EndStartsOnOwnLine$        & FiStartsOnOwnLine         \\
			$\EndFinishesWithLineBreak$  & FiFinishesWithLineBreak   \\
		\end{tabular}
		\\
		\cmidrule{2-3}
		optionalArguments &
		\begin{lstlisting}[escapeinside={(*@}{@*)},nolol=true]
...(*@$\BeginStartsOnOwnLine$@*)
[(*@$\BodyStartsOnOwnLine$@*)
body of opt arg(*@$\EndStartsOnOwnLine$@*)
](*@$\EndFinishesWithLineBreak$@*)
...
  \end{lstlisting}
		&
		\begin{tabular}[t]{c@{~}l@{}}
			$\BeginStartsOnOwnLine$     & LSqBStartsOnOwnLine\footnote{LSqB stands for Left Square Bracket} \\
			$\BodyStartsOnOwnLine$      & OptArgBodyStartsOnOwnLine                                         \\
			$\EndStartsOnOwnLine$       & RSqBStartsOnOwnLine                                               \\
			$\EndFinishesWithLineBreak$ & RSqBFinishesWithLineBreak                                         \\
		\end{tabular}
		\\
		\cmidrule{2-3}
		mandatoryArguments &
		\begin{lstlisting}[escapeinside={(*@}{@*)},nolol=true]
...(*@$\BeginStartsOnOwnLine$@*)
{(*@$\BodyStartsOnOwnLine$@*)
body of mand arg(*@$\EndStartsOnOwnLine$@*)
}(*@$\EndFinishesWithLineBreak$@*)
...
  \end{lstlisting}
		&
		\begin{tabular}[t]{c@{~}l@{}}
			$\BeginStartsOnOwnLine$     & LCuBStartsOnOwnLine\footnote{LCuB stands for Left Curly Brace} \\
			$\BodyStartsOnOwnLine$      & MandArgBodyStartsOnOwnLine                                     \\
			$\EndStartsOnOwnLine$       & RCuBStartsOnOwnLine                                            \\
			$\EndFinishesWithLineBreak$ & RCuBFinishesWithLineBreak                                      \\
		\end{tabular}
		\\
		\cmidrule{2-3}
		commands &
		\begin{lstlisting}[escapeinside={(*@}{@*)},morekeywords={mycommand},nolol=true,]
before words(*@$\BeginStartsOnOwnLine$@*)
\mycommand(*@$\BodyStartsOnOwnLine$@*)
(*@$\langle$\itshape{arguments}$\rangle$@*)
  \end{lstlisting}
		&
		\begin{tabular}[t]{c@{~}l@{}}
			$\BeginStartsOnOwnLine$ & CommandStartsOnOwnLine           \\
			$\BodyStartsOnOwnLine$  & CommandNameFinishesWithLineBreak \\
		\end{tabular}
		\\
		\cmidrule{2-3}
		namedGroupingBraces Brackets &
		\begin{lstlisting}[escapeinside={(*@}{@*)},morekeywords={myname},nolol=true,]
before words(*@$\BeginStartsOnOwnLine$@*)
myname(*@$\BodyStartsOnOwnLine$@*)
(*@$\langle$\itshape{braces/brackets}$\rangle$@*)
  \end{lstlisting}
		&
		\begin{tabular}[t]{c@{~}l@{}}
			$\BeginStartsOnOwnLine$ & NameStartsOnOwnLine       \\
			$\BodyStartsOnOwnLine$  & NameFinishesWithLineBreak \\
		\end{tabular}
		\\
		\cmidrule{2-3}
		keyEqualsValuesBraces &
		\begin{lstlisting}[escapeinside={(*@}{@*)},morekeywords={key},nolol=true,]
before words(*@$\BeginStartsOnOwnLine$@*)
key(*@$\EqualsStartsOnOwnLine$@*)=(*@$\BodyStartsOnOwnLine$@*)
(*@$\langle$\itshape{braces/brackets}$\rangle$@*)
  \end{lstlisting}
		&
		\begin{tabular}[t]{c@{~}l@{}}
			$\BeginStartsOnOwnLine$  & KeyStartsOnOwnLine          \\
			$\EqualsStartsOnOwnLine$ & EqualsStartsOnOwnLine       \\
			$\BodyStartsOnOwnLine$   & EqualsFinishesWithLineBreak \\
		\end{tabular}
		\\
		\cmidrule{2-3}
		items &
		\begin{lstlisting}[escapeinside={(*@}{@*)},nolol=true]
before words(*@$\BeginStartsOnOwnLine$@*)
\item(*@$\BodyStartsOnOwnLine$@*)
...
  \end{lstlisting}
		&
		\begin{tabular}[t]{c@{~}l@{}}
			$\BeginStartsOnOwnLine$ & ItemStartsOnOwnLine       \\
			$\BodyStartsOnOwnLine$  & ItemFinishesWithLineBreak \\
		\end{tabular}
		\\
		\cmidrule{2-3}
		specialBeginEnd &
		\begin{lstlisting}[escapeinside={(*@}{@*)},nolol=true]
before words(*@$\BeginStartsOnOwnLine$@*)
\[(*@$\BodyStartsOnOwnLine$@*)
body of special(*@$\EndStartsOnOwnLine$@*)
\](*@$\EndFinishesWithLineBreak$@*)
after words
  \end{lstlisting}
		&
		\begin{tabular}[t]{c@{~}l@{}}
			$\BeginStartsOnOwnLine$     & SpecialBeginStartsOnOwnLine     \\
			$\BodyStartsOnOwnLine$      & SpecialBodyStartsOnOwnLine      \\
			$\EndStartsOnOwnLine$       & SpecialEndStartsOnOwnLine       \\
			$\EndFinishesWithLineBreak$ & SpecialEndFinishesWithLineBreak \\
		\end{tabular}
		\\
		\bottomrule
	\end{longtable}
